\paragraph{整数温度 pressure 骨架、零点静默、四次 Perron 单位与 \(2\)-进分辨率边界}
在整数阶 moment-kernel 的有限维实现、连续射影压力在整点上的锚定、六倍窗零点密度的精确指数率、Fibonacci 共振整函数的零点自相似、零温冻结端面的四次 Perron 根、\(S_4\)--Hurwitz 链的 Prym 平方化以及分辨率塔的 \(2\)-进逆极限都已固定之后，还可把这些彼此分散的接口继续压缩为下列七条后果。

\begin{theorem}[整数温度样本构成连续压力曲线的代数骨架]\label{thm:xi-time-part9zo-integer-temperature-pressure-algebraic-skeleton}
沿用命题 \ref{prop:pom-moment-kernel-exists} 与定理 \ref{thm:xi-projective-moment-tower-anchors-normalized-pressure} 的记号。
对每个整数 \(q\ge 2\)，记
\[
\mathcal G_q(z):=\sum_{n\ge 0}S_q(m_0+n)z^n.
\]
则有严格恒等式
\[
\mathcal G_q(z)=u_q^{\mathsf T}(I-zA_q)^{-1}v_q\in \QQ(z).
\]
并且对 Fold 主核有整点锚定公式
\[
\Lambda(q-1)=\log r_q-\log 2.
\]
特别地，
\[
2e^{\Lambda(q-1)}=r_q
\]
是一个 Perron 代数整数。换言之，连续压力曲线在全部整数温度点 \(t=q-1\) 上并非任意实值，而是落在一条对数--Perron 代数整数骨架之上。
\end{theorem}
\begin{proof}
由命题 \ref{prop:pom-moment-kernel-exists}，对每个固定 \(q\ge 2\) 都存在有限非负整数矩阵 \(A_q\)、整数向量 \(u_q,v_q\) 与起点 \(m_0\)，使
\[
S_q(m)=u_q^{\mathsf T}A_q^{\,m-m_0}v_q
\qquad (m\ge m_0).
\]
故
\[
\mathcal G_q(z)
=
\sum_{n\ge 0}u_q^{\mathsf T}A_q^n v_q\,z^n
=
u_q^{\mathsf T}\Bigl(\sum_{n\ge 0}(zA_q)^n\Bigr)v_q
=
u_q^{\mathsf T}(I-zA_q)^{-1}v_q.
\]
由于 \(A_q\) 为整数矩阵，\((I-zA_q)^{-1}\) 的各条目都是有理函数，故 \(\mathcal G_q(z)\in \QQ(z)\)。

另一方面，定理 \ref{thm:xi-projective-moment-tower-anchors-normalized-pressure} 给出
\[
\Lambda(q-1)=\log r_q-\log |A_{\mathrm{out}}|.
\]
对 Fold 主核有 \(|A_{\mathrm{out}}|=2\)，因此
\[
\Lambda(q-1)=\log r_q-\log 2.
\]
再由 Perron--Frobenius 定理，\(r_q=\rho(A_q)\) 是整数矩阵 \(A_q\) 的特征值，故为 Perron 代数整数。于是
\[
2e^{\Lambda(q-1)}=r_q
\]
具有所述算术性质。证毕。
\end{proof}

\begin{theorem}[星形 \(q\)-碰撞谱的代数次数至多二次增长]\label{thm:xi-time-part9zo-star-perron-degree-quadratic}
沿用命题 \ref{prop:pom-moment-minrecurrence-hankel}、
定理 \ref{thm:pom-star-moment-kernel-perron-symmetric-compression}
与定理 \ref{thm:pom-sympower-hankel-polynomial-upper-bound} 的记号。
在文稿的星形耦合情形下，记 \(r_q\) 为 \(q\)-重碰撞的 Perron 根，\(d_q\) 为序列 \(\{S_q(m)\}\) 的最小 Hankel 维数。则有显式上界
\[
\deg_{\QQ}(r_q)\le d_q\le \binom{q+1}{2}.
\]
因此，虽然朴素 \(q\)-份同步直积具有指数态空间，但可见谱复杂度与 Perron 数域复杂度都同时塌缩到二次量级。
\end{theorem}
\begin{proof}
由定理 \ref{thm:pom-star-moment-kernel-perron-symmetric-compression}，在星形耦合的三态对称压缩下，\(r_q\) 可由维数
\[
\binom{q+1}{2}
\]
的压缩矩阵 \(\widetilde A_q^{\mathrm{eq}}\) 读出。于是 \(r_q\) 满足其特征多项式，从而
\[
\deg_{\QQ}(r_q)\le \binom{q+1}{2}.
\]

另一方面，命题 \ref{prop:pom-moment-minrecurrence-hankel} 说明 \(d_q\) 等于序列 \(m\mapsto S_q(m)\) 的最小递推阶数，也等于最小 realization 维数；定理 \ref{thm:pom-sympower-hankel-polynomial-upper-bound} 在一般 \(N\)-状态情形下给出
\[
d_q\le \binom{q+N-1}{N-1},
\]
而此处 \(N=3\)，故
\[
d_q\le \binom{q+1}{2}.
\]
再取命题 \ref{prop:pom-moment-minrecurrence-hankel} 给出的最小 realization，其维数为 \(d_q\)。若 \(r_q\) 不是该 realization 矩阵的特征值，则由该 realization 产生的序列只能按严格小于 \(r_q\) 的谱半径增长，这与 \(r_q\) 作为同一碰撞序列的 Perron 主尺度矛盾。故 \(r_q\) 是某个 \(d_q\times d_q\) 整矩阵的特征值，从而
\[
\deg_{\QQ}(r_q)\le d_q.
\]
合并即得
\[
\deg_{\QQ}(r_q)\le d_q\le \binom{q+1}{2}.
\]
证毕。
\end{proof}
\leanverified{paper\_xi\_time\_part9zo\_star\_perron\_degree\_quadratic}

\begin{theorem}[算术零点块具有绝对宏观性而归一化静默]\label{thm:xi-time-part9zo-zero-block-macroscopic-normalized-silent}
沿用推论 \ref{cor:fold-zero-half-index-multiple6}、
定理 \ref{thm:fold-zero-density-sparse}
与定理 \ref{thm:xi-time-part70c-zero-density-exact-ratio-exponent-multiple6} 的记号，并记
\[
\mathcal Z_m:=\{k\in \ZZ/(F_{m+2}\ZZ):\ \widehat d_m(k)=0\}.
\]
记
\[
M:=F_{m+2},
\qquad
\eta_m:=\frac{|\mathcal Z_m|}{F_{m+2}}.
\]
若
\[
m+2\equiv 0\pmod 6,
\]
则有
\[
\eta_m\ge \frac{F_{(m+2)/2}}{F_{m+2}},
\qquad
\eta_m=\tau(m+2)\,O(\varphi^{-m/2}),
\]
从而
\[
\log(M-|\mathcal Z_m|)
=
\log M+O\!\bigl(\tau(m+2)\varphi^{-m/2}\bigr),
\]
以及
\[
\sqrt{M-1-|\mathcal Z_m|}
=
\sqrt M\Bigl(1+O\!\bigl(\tau(m+2)\varphi^{-m/2}\bigr)\Bigr).
\]
因此，六倍窗上虽出现规模达 \(\asymp \varphi^{m/2}\) 的零点大陪集块，但凡只经由 \(|\mathcal Z_m|\) 或其归一化比例进入的主阶估计，都只能获得半指数级副修正。
\end{theorem}
\begin{proof}
下界由推论 \ref{cor:fold-zero-half-index-multiple6} 直接给出：
\[
|\mathcal Z_m|\ge F_{(m+2)/2},
\]
故
\[
\eta_m\ge \frac{F_{(m+2)/2}}{F_{m+2}}.
\]
另一方面，定理 \ref{thm:fold-zero-density-sparse} 给出
\[
\eta_m=\frac{|\mathcal Z_m|}{F_{m+2}}
\le
\tau(m+2)\frac{F_{\lfloor (m+2)/2\rfloor}}{F_{m+2}}
=
\tau(m+2)\,O(\varphi^{-m/2}),
\]
而定理 \ref{thm:xi-time-part70c-zero-density-exact-ratio-exponent-multiple6} 进一步表明该指数率恰为 \(-\frac12\log\varphi\)。

于是
\[
M-|\mathcal Z_m|=M(1-\eta_m),
\]
从而
\[
\log(M-|\mathcal Z_m|)
=
\log M+\log(1-\eta_m)
=
\log M+O(\eta_m).
\]
再注意
\[
M-1-|\mathcal Z_m|
=
M\Bigl(1-\eta_m-\frac1M\Bigr),
\]
且 \(M^{-1}=O(\varphi^{-m})=O(\eta_m)\)，于是
\[
\sqrt{M-1-|\mathcal Z_m|}
=
\sqrt M\sqrt{1-\eta_m-\frac1M}
=
\sqrt M\bigl(1+O(\eta_m)\bigr).
\]
把 \(\eta_m=\tau(m+2)\,O(\varphi^{-m/2})\) 代回即得两式。最后一句只是指出：凡通过 \(M-|\mathcal Z_m|\)、\(\sqrt{M-1-|\mathcal Z_m|}\) 或等价归一化比例进入的零点驱动上界，其主阶都无法被 \(|\mathcal Z_m|\) 改写。证毕。
\end{proof}

\begin{theorem}[Fibonacci 共振零点满足精确重整化与 Weyl 型计数律]\label{thm:xi-time-part9zo-fibonacci-cosine-zeros-weyl-law}
沿用定理 \ref{thm:xi-fold-resonance-zero-count-selfsimilar-linear} 的记号，记
\[
\mathcal F_\varphi(z)=\prod_{n=2}^{\infty}\cos(\pi z\varphi^{-n}),
\qquad
N_\varphi^+(R):=\#\bigl\{x\in \mathcal Z(\mathcal F_\varphi):\ 0<x\le R\bigr\}.
\]
则对每个 \(R>0\) 都有严格递推
\[
N_\varphi^+(R)=N_\varphi^+\!\left(\frac{R}{\varphi}\right)+\left\lfloor \frac{R}{\varphi^2}+\frac12\right\rfloor.
\]
并且
\[
N_\varphi^+(R)=R+O(\log R),
\qquad
N_\varphi(R):=\#\bigl(\mathcal Z(\mathcal F_\varphi)\cap[-R,R]\bigr)=2R+O(\log R).
\]
换言之，Fibonacci 共振零点并非给出稀薄谱，而是具有严格线性密度；其主斜率恰由
\[
\sum_{n\ge 2}\varphi^{-n}=1
\]
唯一锁定。
\end{theorem}
\begin{proof}
定理 \ref{thm:xi-fold-resonance-zero-count-selfsimilar-linear} 已经给出精确自相似递推
\[
N_\varphi^+(R)=N_\varphi^+\!\left(\frac{R}{\varphi}\right)+\left\lfloor \frac{R}{\varphi^2}+\frac12\right\rfloor,
\]
以及渐近式
\[
N_\varphi^+(R)=R+O(\log R),
\qquad
N_\varphi(R)=2R+O(\log R).
\]
另一方面，由同一定理中的精确求和公式
\[
N_\varphi^+(R)=
\sum_{n=2}^{\lfloor \log_\varphi(2R)\rfloor}
\left\lfloor R\varphi^{-n}+\frac12\right\rfloor,
\]
把整部份替换为 \(R\varphi^{-n}+O(1)\) 并注意有效层数只有 \(O(\log R)\)，可得
\[
N_\varphi^+(R)
=
R\sum_{n\ge 2}\varphi^{-n}+O(\log R).
\]
而
\[
\sum_{n\ge 2}\varphi^{-n}=1,
\]
故线性主项恰为 \(R\)。对称性再给出
\[
N_\varphi(R)=2N_\varphi^+(R)=2R+O(\log R).
\]
证毕。
\end{proof}
\leanverified{paper\_xi\_time\_part9zo\_fibonacci\_cosine\_zeros\_weyl\_law}

\begin{theorem}[零温饱和端点常数是四次 Perron 单位的平方根]\label{thm:xi-time-part9zo-zero-temp-quartic-perron-unit}
沿用定理 \ref{thm:xi-terminal-real-input-positive-entropy-freezing-face} 的记号，令
\[
\lambda_\infty:=\rho(B_\infty)=c^2.
\]
则有
\[
\lambda_\infty^4-6\lambda_\infty^3+9\lambda_\infty^2-\lambda_\infty-1=0,
\]
从而
\[
c^8-6c^6+9c^4-c^2-1=0.
\]
特别地，\(\lambda_\infty\) 是一个 Perron 代数单位；因此零温饱和端点常数 \(c\) 并非单纯数值参数，而是四状态地基 SFT 的平方根 Perron 单位。
\end{theorem}
\begin{proof}
定理 \ref{thm:xi-terminal-real-input-positive-entropy-freezing-face} 已给出
\[
\lambda_\infty=\rho(B_\infty)=y,
\qquad
y^4-6y^3+9y^2-y-1=0,
\qquad
c=\sqrt y.
\]
把 \(y=\lambda_\infty\) 代回即得
\[
\lambda_\infty^4-6\lambda_\infty^3+9\lambda_\infty^2-\lambda_\infty-1=0.
\]
再以 \(\lambda_\infty=c^2\) 代入，得到
\[
c^8-6c^6+9c^4-c^2-1=0.
\]

由于
\[
x^4-6x^3+9x^2-x-1
\]
是首一整系数多项式且常数项为 \(-1\)，\(\lambda_\infty\) 的最小多项式作为其首一整系数因子，常数项必为 \(\pm 1\)。因此 \(\lambda_\infty\) 是代数单位。又因为 \(\lambda_\infty=\rho(B_\infty)\) 来自非负整数矩阵 \(B_\infty\) 的谱半径，故它同时是一个 Perron 数。证毕。
\end{proof}

\begin{theorem}[\texorpdfstring{$S_4$}{S4}--Hurwitz 链上的三重 Prym 不可能绝对单]\label{thm:xi-time-part9zo-prym-c6h-not-absolutely-simple}
沿用结论 \ref{con:xi-time-part9n-s4-all-subgroup-genera} 与定理 \ref{thm:killo-s4-burnside-kani-rosen-prym-square} 的记号，令
\[
P:=\operatorname{Prym}(C_6/H).
\]
则存在亏格 \(4\) 曲线 \(Z=X/D_8\)，使得
\[
P\sim J(Z)^2.
\]
特别地，
\[
\dim P=8,
\qquad
\End^0(P)\supset M_2(\QQ),
\]
从而 \(P\) 不可能绝对单。
\end{theorem}
\begin{proof}
由定理 \ref{thm:killo-s4-burnside-kani-rosen-prym-square}，
\[
\operatorname{Prym}(C_6/H)\sim J(X/D_8)^2.
\]
又由结论 \ref{con:xi-time-part9n-s4-all-subgroup-genera}，
\[
g(X/D_8)=4.
\]
记
\[
Z:=X/D_8,
\]
便得到
\[
P\sim J(Z)^2
\]
且
\[
\dim P=2g(Z)=8.
\]

对任意阿贝尔簇 \(A\)，都有自然嵌入
\[
M_2(\QQ)\hookrightarrow \End^0(A^2).
\]
将 \(A\) 取为 \(J(Z)\)，再利用同源不改变有理自同态代数的半单矩阵型包含，便得
\[
\End^0(P)\supset M_2(\QQ).
\]
若 \(P\) 绝对单，则 \(\End^0(P)\) 不可能含有非平凡矩阵代数 \(M_2(\QQ)\)。故 \(P\) 不绝对单。证毕。
\end{proof}

\begin{theorem}[分辨率边界的柱集代数与 \(2\)-进闭球代数完全一致]\label{thm:xi-time-part9zo-dyadic-boundary-cylinder-ball-identification}
\leanverified{paper\_xi\_time\_part9zo\_dyadic\_boundary\_cylinder\_ball\_identification}
记
\[
\partial Q:=\varprojlim_n Q_{2n}
\]
为超立方体图塔在截断映射下的逆极限。则在相容基选择下有自然同构
\[
\partial Q\cong T_2(E_{\min})\cong \ZZ_2^2.
\]
更强地，固定前 \(n\) 层二进前缀所得到的柱集，恰对应于 \(\ZZ_2^2\) 中某个模 \(2^n\) 剩余类
\[
(a,b)+2^n\ZZ_2^2,
\]
亦即最大 \(2\)-进范数下半径 \(2^{-n}\) 的一个闭球。因此，分辨率边界上的柱集代数与 \(\ZZ_2^2\) 的闭球代数完全一致。
\end{theorem}
\begin{proof}
推论 \ref{cor:killo-2adic-tate-hypercube-limit} 已经给出，在相容基
\[
P_n=2P_{n+1},
\qquad
Q_n=2Q_{n+1}
\]
之下，
\[
\varprojlim_n E_{\min}[2^n]\cong T_2(E_{\min})\cong \ZZ_2^2,
\]
而同一相容识别把图塔 \((Q_{2n})_{n\ge 1}\) 的逆极限送到该 \(2\)-进地址空间。故
\[
\partial Q\cong T_2(E_{\min})\cong \ZZ_2^2.
\]

现在固定前 \(n\) 层前缀。按照推论 \ref{cor:killo-2adic-tate-hypercube-limit} 的坐标解释，这等价于固定
\[
(a\bmod 2^n,\ b\bmod 2^n)\in (\ZZ/2^n\ZZ)^2.
\]
因此对应子集正是
\[
(a,b)+2^n\ZZ_2^2.
\]
在 \(\ZZ_2^2\) 的最大 \(2\)-进范数
\[
\|(x,y)\|_{\max}:=\max\{|x|_2,|y|_2\}
\]
下，这恰是以任一代表点为中心、半径 \(2^{-n}\) 的闭球。于是柱集的有限交并补仍与闭球余类结构一致，生成的代数正是 \(2\)-进闭球代数。证毕。
\end{proof}

\begin{theorem}[零点块对持久大纤维偏置的无力性]\label{thm:xi-time-part9zoa-zero-block-maxfiber-impotence}
\leanverified{paper\_xi\_time\_part9zoa\_zero\_block\_maxfiber\_impotence}
沿用定理 \ref{thm:xi-time-part9zo-zero-block-macroscopic-normalized-silent}、
定理 \ref{thm:xi-time-part60ab-window6-zero-block-halfscale-information-improvement}
与定理 \ref{thm:xi-time-part9m-maxfiber-uniformity-gap-chain} 的记号。若
\[
M:=F_{m+2},
\qquad
\mathcal Z_m:=\{k\in \ZZ/(M\ZZ):\ \widehat d_m(k)=0\},
\qquad
\eta_m:=\frac{|\mathcal Z_m|}{M},
\]
并且
\[
m+2\equiv 0\pmod 6,
\]
则虽有确定性的半阶零点块下界
\[
|\mathcal Z_m|\ge F_{(m+2)/2},
\]
但仍满足
\[
\eta_m=O(\varphi^{-m/2}).
\]
更具体地，任何仅通过
\[
M-|\mathcal Z_m|,
\qquad
\log(M-|\mathcal Z_m|),
\qquad
\sqrt{M-1-|\mathcal Z_m|}
\]
进入的零点驱动上界，其相对于零点自由基线的改变量都至多具有 \(O(\eta_m)=O(\varphi^{-m/2})\) 的相对尺度。与此相对，最大纤维偏置满足
\[
\liminf_{m\to\infty}\frac{M_m}{\overline d_m}\ge 1+C_\varphi>1.
\]
因此，算术零点块无法消除持久的 \(L^\infty\) 级非均匀性；真正顽固的障碍并不在零点集合，而在大纤维偏置本身。
\end{theorem}
\begin{proof}
由定理 \ref{thm:xi-time-part9zo-zero-block-macroscopic-normalized-silent}，
\[
\eta_m=\frac{|\mathcal Z_m|}{M}
=
\tau(m+2)\,O(\varphi^{-m/2})
=
O(\varphi^{-m/2}),
\]
其中最后一步仅用到约数函数的次幂增长远慢于任意指数函数。

同一定理还给出
\[
M-|\mathcal Z_m|=M(1-\eta_m),
\]
\[
\log(M-|\mathcal Z_m|)=\log M+O(\eta_m),
\]
\[
\sqrt{M-1-|\mathcal Z_m|}
=
\sqrt M\bigl(1+O(\eta_m)\bigr).
\]
故凡是仅通过这三类表达式进入的 zero-driven 上界，其与零点自由主项相比都只能产生 \(O(\eta_m)\) 级的相对改动。这正是“零点块只能提供半指数级副修正”的精确形式。

另一方面，定理 \ref{thm:xi-time-part9m-maxfiber-uniformity-gap-chain} 已经给出
\[
\liminf_{m\to\infty}\frac{M_m}{\overline d_m}\ge 1+C_\varphi>1.
\]
因此最大纤维相对均值的偏置在极限中仍与 \(1\) 保持严格正距离，不可能被任何 \(O(\varphi^{-m/2})\) 级的 zero-block 修正所抹平。于是折叠分布的持久非均匀性并非由零点块驱动，真正的主障碍只能来自大纤维偏置。证毕。
\end{proof}

\endinput


\endinput
